% ============================================================
%  Section 1: Business Intelligence & Data Warehouse
% ============================================================
\section{Business Intelligence \& Data Warehouse}

% -------- subsection: What is BI? --------

\begin{frame}
  \frametitle{What is Business Intelligence?}

  \begin{block}{Definition}
    \textbf{Business Intelligence (BI)} refers to all tools and applications with
    a \textbf{decision-supporting character} that lead to better insight into
    one's own business and a deeper understanding of the mechanisms in relevant
    chains of causality.
    \hfill\small(Weber et al., 1999)
  \end{block}

  \vspace{0.4em}
  \begin{block}{In plain terms}
    BI is the \textbf{intelligence service of the company}. It is the process
    that generates knowledge from fragmented, heterogeneous corporate, market,
    and competitive data about one's own and others' positions, potentials,
    and perspectives.
  \end{block}
\end{frame}

% ------------------------------------------------------------
\begin{frame}
  \frametitle{Two Ways to Think About BI}

  \begin{block}{Wide BI Understanding (Technology \& Application)}
    \begin{itemize}
      \item ETL, Data Warehouse, OLAP, Reporting
      \item Standard \& ad-hoc reporting, EIS/DSS/MIS systems
      \item Data Mining, Text Mining
      \item Planning, Consolidation, Balanced Scorecard
      \item Analytical CRM
    \end{itemize}
  \end{block}

  \begin{block}{Narrow BI Understanding (Analytics Focus)}
    \begin{itemize}
      \item Focused on KPI systems and decision support
      \item Management reporting and scorecards
      \item Analytical applications for specific business domains
    \end{itemize}
  \end{block}

  % PICTURE PLACEHOLDER
  \begin{alertblock}{Picture: BI Facets Diagram}
    \textit{Add the ``Facets of Business Intelligence'' diagram showing the
    phases (data provision, analytics) on vertical axis and the spectrum from
    narrow to wide BI on horizontal axis. Include ETL, DWH, Reporting, OLAP,
    Data Mining, BSC positions.}
  \end{alertblock}
\end{frame}

% -------- subsection: Why BI? --------

\begin{frame}
  \frametitle{Why Do Companies Need BI?}

  \begin{columns}[T]
    \column{0.5\textwidth}
    \begin{block}{Common Complaints}
      \begin{itemize}
        \item Reports arrive \textbf{too late}
        \item Numbers are \textbf{inconsistent} across departments
        \item Data is \textbf{inaccurate} or incomplete
        \item Analysis is done manually in Excel --- error-prone and slow
        \item Every department has its own \textit{version of the truth}
      \end{itemize}
    \end{block}

    \column{0.5\textwidth}
    \begin{block}{The Finance Reporting Problem}
      Without a DWH, a controller spends up to \textbf{48\,\%} of their time
      on data collection and validation --- instead of on actual analysis.
      \hfill\small(Hackett Group Study)
    \end{block}

    \begin{alertblock}{Goal of BI}
      Fast, consistent, fact-based information for everyone who needs to
      make decisions.
    \end{alertblock}
  \end{columns}
\end{frame}

% ------------------------------------------------------------
\begin{frame}
  \frametitle{Information Supply Chain}

  Just as goods flow through a supply chain, \textbf{information flows} through
  the organisation.

  % PICTURE PLACEHOLDER
  \begin{alertblock}{Picture: Information Supply Chain}
    \textit{Add the Information Supply Chain diagram:
    Data $\to$ Analysis $\to$ Knowledge $\to$ Action, with ``Operative Business
    Processes'' underneath as the source. Show the cycle.}
  \end{alertblock}

  \begin{block}{Roles in the Information Supply Chain}
    \begin{tabularx}{\linewidth}{@{}l X@{}}
      \toprule
      \textbf{Role}          & \textbf{Responsibility} \\
      \midrule
      Business User          & Consumes business information \\
      Reporter               & Creates standard reports \\
      Guide                  & Provides detailed information \\
      Miner / Analyst        & Data analysis and discovery \\
      Administrator          & Infrastructure and application management \\
      \bottomrule
    \end{tabularx}
  \end{block}
\end{frame}

% -------- subsection: OLTP brief note --------

\begin{frame}
  \frametitle{Where Does the Data Come From? Operational Systems (OLTP)}

  \begin{block}{Online Transaction Processing (OLTP)}
    The \textbf{source} of data for the DWH. OLTP systems run the day-to-day
    business: orders, payments, inventory, HR.
    They are optimised for \textbf{fast writes and concurrent users}, not for analytics.
  \end{block}

  \begin{tabularx}{\linewidth}{@{}l X X@{}}
    \toprule
    \textbf{Property}  & \textbf{OLTP (Source)}            & \textbf{DWH / OLAP (Target)} \\
    \midrule
    Purpose            & Run the business                  & Analyse the business \\
    Operations         & INSERT/UPDATE/DELETE              & Mainly SELECT \\
    Data               & Current, single record            & Historical, aggregated \\
    Volume             & MB -- GB                          & GB -- TB \\
    Users              & Clerks, ERP apps                  & Analysts, managers \\
    Response time      & Milliseconds                      & Seconds -- minutes \\
    Examples           & SAP, Salesforce, POS              & DWH, OLAP cubes, BI tools \\
    \bottomrule
  \end{tabularx}
\end{frame}

% -------- subsection: DWH Definition --------

\begin{frame}
  \frametitle{The Data Warehouse: Definition}

  \begin{block}{Bill Inmon's Definition (1992)}
    A Data Warehouse is a \textbf{subject-oriented}, \textbf{integrated},
    \textbf{time-variant}, and \textbf{non-volatile} collection of data
    in support of management's decision-making process.
  \end{block}

  \vspace{0.5em}
  \begin{block}{In practice, a DWH is\ldots}
    \begin{itemize}
      \item A central \textbf{data pool} that forms the basis of information
        supply in the enterprise
      \item Sourced from operational systems (ERP like SAP) and external
        data (market data, competitor information)
      \item The \textbf{foundation} from which reports, analyses, and dashboards
        are built
    \end{itemize}
  \end{block}
\end{frame}

% ------------------------------------------------------------
\begin{frame}
  \frametitle{The Four Characteristics of a DWH}

  \begin{block}{\textbf{1. Subject-Oriented}}
    Data is organised around key business subjects --- Customer, Product,
    Sales, Service --- not around applications or processes.
    Operational systems are process-oriented; the DWH is subject-oriented.
  \end{block}

  \begin{block}{\textbf{2. Integrated}}
    Data from heterogeneous sources is cleaned, transformed, and stored
    in a \textbf{consistent, unified format}.
    Names, units, and encodings are harmonised across all sources.
  \end{block}

  % PICTURE PLACEHOLDER
  \begin{alertblock}{Picture: Subject-Oriented vs.\ Operational}
    \textit{Add a diagram showing: Operational apps (Sales, Finance,
    Marketing, Hotline, Training) on the left pointing into the DWH
    on the right, which is organised by subjects: Customer, Service, Product.}
  \end{alertblock}
\end{frame}

% ------------------------------------------------------------
\begin{frame}
  \frametitle{The Four Characteristics of a DWH (cont.)}

  \begin{block}{\textbf{3. Time-Variant}}
    Every data record is associated with a \textbf{time stamp or validity period}.
    The DWH stores \textbf{multiple snapshots} of data over time --- unlike
    operational systems, which typically reflect only the current state.
    Years or even decades of history are retained.
  \end{block}

  \begin{block}{\textbf{4. Non-Volatile}}
    Data is \textbf{loaded} into the DWH and then \textbf{only read}.
    There are no updates or deletes in the transactional sense.
    The data is stable and persistent within a time unit.
    Loading happens in \textbf{defined intervals} (daily, weekly, \ldots).
  \end{block}

  \begin{examples}{Summary}
    The DWH stores \textit{detailed, filtered, subject-related data},
    derived aggregates and reports, metadata, historical data (multiple years),
    and pre-aggregated data for fast querying.
  \end{examples}
\end{frame}

% -------- subsection: DWH Application Areas --------

\begin{frame}
  \frametitle{Application Areas of the Data Warehouse}

  The most common use cases are in the business / management area.

  \begin{block}{Main Areas}
    \begin{itemize}
      \item \textbf{Controlling:} Data basis for financial reporting, KPIs,
        budget vs.\ actuals, consolidation
      \item \textbf{Customer Relationship Management (CRM):}
        Customer segmentation, direct marketing, churn prediction,
        customer profitability analysis
      \item \textbf{Knowledge Management:}
        Combining structured data (DWH) with unstructured data
        (documents, email, images) via document management and intranets
      \item \textbf{Management Concepts:}
        Balanced Scorecard, Business Performance Management,
        Total Quality Management support
    \end{itemize}
  \end{block}
\end{frame}

% ------------------------------------------------------------
\begin{frame}
  \frametitle{Use Case: Controlling}

  \begin{block}{Typical Pain Points in Controlling without a DWH}
    \begin{itemize}
      \item Heterogeneous data sources: ERP, SAP, Excel, local databases
      \item Validity of financial figures is questionable
      \item Manual consolidation of data from many systems
      \item High time expenditure for report preparation
      \item Up to \textbf{48\,\%} of controller time goes to data gathering,
        not analysis \hfill\small(Hackett Group)
    \end{itemize}
  \end{block}

  \begin{alertblock}{Goal of Controlling in BI}
    The DWH automates data consolidation and provides a \textbf{single
    source of truth} so controllers can focus on actual analysis and
    decision support.
  \end{alertblock}
\end{frame}

% ------------------------------------------------------------
\begin{frame}
  \frametitle{Use Case: Customer Relationship Management (CRM)}

  \begin{block}{The CRM Data Challenge}
    Customer data is scattered across many systems:
    \begin{itemize}
      \item Transactional data in operational applications (orders, accounts)
      \item Offers in Word documents
      \item Revenue / product statistics in the DWH
      \item Contact notes as emails and private files
      \item Customer self-presentation on the web
    \end{itemize}
  \end{block}

  \begin{block}{What BI enables in CRM}
    \begin{itemize}
      \item Customer segmentation and profiling
      \item Direct marketing campaign management
      \item Shopping basket / cross-sell analysis
      \item Customer profitability and credit scoring
      \item Churn (defection) prediction
      \item Service quality and response analysis
    \end{itemize}
  \end{block}
\end{frame}

% -------- subsection: ROI --------

\begin{frame}
  \frametitle{ROI and Cost Considerations}

  \begin{block}{Return on Investment (ROI)}
    \begin{itemize}
      \item A measure of the profitability of invested capital
      \item Only quantifiable aspects are considered
      \item DWH projects often have a difficult-to-quantify benefit side
      \item Expected ROI materialises \textbf{over time} --- not immediately
    \end{itemize}
  \end{block}

  % PICTURE PLACEHOLDER
  \begin{alertblock}{Picture: ROI Curve over Time}
    \textit{Add a cost/benefit curve diagram: X-axis = time, Y-axis = cost/benefit.
    Show the initial investment dip, the break-even point, and the rising
    return-on-investment as the DWH matures.}
  \end{alertblock}

  \begin{block}{Typical Cost Categories}
    \begin{itemize}
      \item \textbf{Up-front:} Planning, modelling, design
      \item \textbf{One-time:} Hardware, software, data cleansing
      \item \textbf{Ongoing:} Data maintenance, data quality, DWH management,
        training, support
    \end{itemize}
  \end{block}
\end{frame}
